\section{Position Representation: por qué la self-attention debe recuperar el orden}

La self-attention es permutation equivariant respecto a la entrada: si se permutan a la vez los tokens y las posiciones, la attention pura no puede conocer el orden original. El lenguaje, la música y las imágenes dependen de la posición relativa, por lo que el modelo debe inyectar una position signal. El encoding learned/sinusoidal del Transformer original es el punto de partida; después, el relative embedding, el relative bias y la rotation escriben la distancia directamente en la interaction.

\subsection{Absolute Position: Learned y Sinusoidal}

El método original suma un position vector $p_i$ al token embedding $x_i$. Un grupo de dimensiones del sinusoidal encoding es:
\[
p_{i,2k}=\sin\!\left(i/10000^{2k/d}\right),\qquad
p_{i,2k+1}=\cos\!\left(i/10000^{2k/d}\right).
\]
Donde $i$ es la posición, $k$ es el índice de frecuencia y $d$ es el model width. No requiere aprendizaje y puede calcularse para cualquier posición, pero que el model realmente extrapole más allá de la longitud de entrenamiento sigue dependiendo del attention behavior y de la distribución de entrenamiento.

\lecturefigure{slide-27-evolution-position-encodings.jpg}{Permutation invariance y el position encoding learned/sinusoidal original}{Diapositiva V027 recuperada de la grabación oficial; 00:33:45}

\begin{knowledgebox}{Lectura de la figura: la position signal debe responder «a qué distancia relativa», no solo «en qué posición absoluta»}
La arquitectura original, a la derecha, agrega el position encoding en la parte inferior. Primero observe si el index absoluto puede expresar de forma natural $i-j$; después, si hay una estructura estable más allá de la longitud de entrenamiento. El learned embedding queda fácilmente limitado por la max length; la sinusoid tiene forma analítica, pero no garantiza que el modelo aprenda una extrapolation confiable.
\end{knowledgebox}

\teachervoice{El ponente señala que la position sinusoidal/learned original no captura bien la relative distance ni la extrapolation, y que esto se volvió el foco de la evolución posterior. Ajustar la base frequency puede cambiar el comportamiento ante la longitud, pero el position scheme no es una perilla mágica independiente: debe coordinarse con el context de entrenamiento y la attention scale.}

\subsection{Relative Embedding, Bias y Rotation}

Los métodos relative incorporan la distancia $r=i-j$ directamente en el score o en el value. Una relative-key attention simplificada es:
\[
s_{ij}=q_i^\top(k_j+a_{i-j}),
\]
donde $a_{i-j}$ es el embedding de la distancia relativa. El relative bias suma un escalar directamente a $s_{ij}$, mientras que el rotary position embedding (RoPE) rota query/key según la posición, de modo que el producto interno contenga la phase relativa; los tres hacen concesiones distintas en parámetros, extrapolación e implementación.

\lecturefigure{slide-28-position-encoding-variants.jpg}{Relative embedding, relative bias y relative rotation}{Diapositiva V028 recuperada de la grabación oficial; 00:36:45}

\begin{knowledgebox}{Términos clave: diferencias de interfaz entre los esquemas de position}
\begin{tabularx}{\textwidth}{@{}L{0.20\textwidth}X X@{}}
\toprule
Esquema & Dónde se coloca la posición & Concesión principal \\
\midrule
Absolute embedding & Se suma al hidden state del token & Simple, pero la distancia se aprende de forma indirecta \\
Relative embedding/bias & Se incorpora al attention score/value pairwise & Expresa la distancia directamente; hay que manejar el rango de distancias \\
Rotary encoding & Rota el subspace de query/key & El producto interno contiene de forma natural la phase relativa; la extrapolación depende del diseño de frecuencias \\
\bottomrule
\end{tabularx}
\end{knowledgebox}

\subsection{Music Transformer: cómo la posición relativa mejora la estructura larga}

La generación musical exige que el modelo mantenga el ritmo, los motivos y las estructuras repetidas; las dependencias entre tokens pueden abarcar lapsos muy largos. Music Transformer usa relative attention, lo que facilita al modelo leer la historia según intervalos y relaciones de repetición; muestra que el position mechanism puede cambiar la qualitative coherence, y no solo las cifras de benchmark.

\lecturefigure{slide-29-music-transformer.jpg}{Music Transformer modela la estructura musical de largo alcance con relative attention}{Diapositiva V029 recuperada de la grabación oficial; 00:37:12}

\begin{knowledgebox}{Lectura de la figura: primero la relation, después escuchar/ver el outcome}
La figura destaca la score decomposition de la relative attention. Primero siga el intervalo temporal entre la note actual y las notes históricas; después compárelo con la vanilla attention. La relative signal aporta un prior estructural, pero la calidad musical final también depende de la tokenización, el training corpus, el sampling y el método de evaluación.
\end{knowledgebox}

\lecturefigure{slide-30-music-transformer-demo.jpg}{Página de qualitative samples de Music Transformer y otros modelos}{Diapositiva V030 recuperada de la grabación oficial; 00:37:18}

\begin{warningbox}{El demo es evidencia, pero de fuerza limitada}
Los ejemplos de la página web ayudan al público a percibir la repetition, la coherence y el drift, pero son susceptibles al cherry-picking, al fragmento reproducido y a las preferencias subjetivas. Conviene combinarlos con la held-out likelihood, métricas estructurales, un blind listening study y ejemplos con varias seeds, en lugar de tomar una pieza musical agradable como una victoria general.
\end{warningbox}

\teachervoice{El ponente reproduce/muestra samples musicales para que el efecto de la relative position sea perceptible; no trata el demo como una metric completa. La clase conserva esta evidencia cualitativa, pero debe quedar claro que solo permite plantear la observación de que «la estructura de largo alcance es mejor», y no basta para determinar una causa única.}

\subsection{Resumen de la sección}

La self-attention por sí misma no conoce el orden. El absolute encoding permite al modelo ver el index; relative/bias/rotation hacen que la pairwise interaction perciba directamente la distancia, y Music Transformer ofrece una aplicación a estructuras largas. La position representation es un componente central de la attention interface, no una suma decorativa que se hace una sola vez antes de la entrada.
